\documentclass[10pt,a4paper]{amsart}
\usepackage{amsmath,amssymb,amsthm,graphicx}
% For commutative diagrams we use Xy-pic
\usepackage[all]{xy}

% We use 2cell for 2-commutative diagrams.
\xyoption{2cell}
% Theorem environments.
%
\theoremstyle{plain}
\newtheorem{theorem}[subsection]{Theorem}
\newtheorem{proposition}[subsection]{Proposition}
\newtheorem{lemma}[subsection]{Lemma}

\theoremstyle{definition}
\newtheorem{definition}[subsection]{Definition}
\newtheorem{example}[subsection]{Example}
\newtheorem{exercise}[subsection]{Exercise}
\newtheorem{situation}[subsection]{Situation}

\theoremstyle{remark}
\newtheorem{remark}[subsection]{Remark}
\newtheorem{remarks}[subsection]{Remarks}

\numberwithin{equation}{subsection}

% Macros
%
\def\lim{\mathop{\mathrm{lim}}\nolimits}
\def\colim{\mathop{\mathrm{colim}}\nolimits}
\def\Spec{\mathop{\mathrm{Spec}}}
\def\Hom{\mathop{\mathrm{Hom}}\nolimits}
\def\Ext{\mathop{\mathrm{Ext}}\nolimits}
\def\SheafHom{\mathop{\mathcal{H}\!\mathit{om}}\nolimits}
\def\SheafExt{\mathop{\mathcal{E}\!\mathit{xt}}\nolimits}
\def\Sch{\mathit{Sch}}
\def\Mor{\mathop{\mathrm{Mor}}\nolimits}
\def\Ob{\mathop{\mathrm{Ob}}\nolimits}
\def\Sh{\mathop{\mathit{Sh}}\nolimits}
\def\NL{\mathop{N\!L}\nolimits}
\def\CH{\mathop{\mathrm{CH}}\nolimits}
\def\proetale{{pro\text{-}\acute{e}tale}}
\def\etale{{\acute{e}tale}}
\def\QCoh{\mathit{QCoh}}
\def\Ker{\mathop{\mathrm{Ker}}}
\def\Im{\mathop{\mathrm{Im}}}
\def\Coker{\mathop{\mathrm{Coker}}}
\def\Coim{\mathop{\mathrm{Coim}}}

% Boxtimes
%
\DeclareMathSymbol{\boxtimes}{\mathbin}{AMSa}{"02}

%
% Macros for moduli stacks/spaces
%
\def\QCohstack{\mathcal{QC}\!\mathit{oh}}
\def\Cohstack{\mathcal{C}\!\mathit{oh}}
\def\Spacesstack{\mathcal{S}\!\mathit{paces}}
\def\Quotfunctor{\mathrm{Quot}}
\def\Hilbfunctor{\mathrm{Hilb}}
\def\Curvesstack{\mathcal{C}\!\mathit{urves}}
\def\Polarizedstack{\mathcal{P}\!\mathit{olarized}}
\def\Complexesstack{\mathcal{C}\!\mathit{omplexes}}
% \Pic is the operator that assigns to X its picard group, usage \Pic(X)
% \Picardstack_{X/B} denotes the Picard stack of X over B
% \Picardfunctor_{X/B} denotes the Picard functor of X over B
\def\Pic{\mathop{\mathrm{Pic}}\nolimits}
\def\Picardstack{\mathcal{P}\!\mathit{ic}}
\def\Picardfunctor{\mathrm{Pic}}
\def\Deformationcategory{\mathcal{D}\!\mathit{ef}}
\usepackage{cleveref}
\crefname{originallemma}{original lemma}{original lemmas}
\Crefname{originallemma}{Original lemma}{Original lemmas}
\makeatletter
\newlabel{homology-lemma-characterize-weak-serre-subcategory}{{0754}{source}{}{sourcepointer-1}{}}
\newlabel{homology-lemma-characterize-weak-serre-subcategory@cref}{{[originallemma][0][]0754}{[1][0][]source}{}{}{}}
\makeatother
\crefname{originalreference}{original reference}{original references}
\Crefname{originalreference}{Original reference}{Original references}
\makeatletter
\newlabel{lemma-cohomological-descent-etale}{{0DDI}{source}{}{sourcepointer-2}{}}
\newlabel{lemma-cohomological-descent-etale@cref}{{[originalreference][0][]0DDI}{[1][0][]source}{}{}{}}
\makeatother
\crefname{originalreference}{original reference}{original references}
\Crefname{originalreference}{Original reference}{Original references}
\makeatletter
\newlabel{lemma-compare-higher-direct-image-proper}{{0DDJ}{source}{}{sourcepointer-3}{}}
\newlabel{lemma-compare-higher-direct-image-proper@cref}{{[originalreference][0][]0DDJ}{[1][0][]source}{}{}{}}
\makeatother
\crefname{originalreference}{original reference}{original references}
\Crefname{originalreference}{Original reference}{Original references}
\makeatletter
\newlabel{lemma-describe-pullback-pi-fppf}{{0DDL}{source}{}{sourcepointer-4}{}}
\newlabel{lemma-describe-pullback-pi-fppf@cref}{{[originalreference][0][]0DDL}{[1][0][]source}{}{}{}}
\makeatother
\crefname{originallemma}{original lemma}{original lemmas}
\Crefname{originallemma}{Original lemma}{Original lemmas}
\makeatletter
\newlabel{more-morphisms-lemma-dominate-fppf-etale-locally}{{0DET}{source}{}{sourcepointer-5}{}}
\newlabel{more-morphisms-lemma-dominate-fppf-etale-locally@cref}{{[originallemma][0][]0DET}{[1][0][]source}{}{}{}}
\makeatother
\crefname{originallemma}{original lemma}{original lemmas}
\Crefname{originallemma}{Original lemma}{Original lemmas}
\makeatletter
\newlabel{morphisms-lemma-base-change-finite}{{01WL}{source}{}{sourcepointer-6}{}}
\newlabel{morphisms-lemma-base-change-finite@cref}{{[originallemma][0][]01WL}{[1][0][]source}{}{}{}}
\makeatother
\crefname{originalreference}{original reference}{original references}
\Crefname{originalreference}{Original reference}{Original references}
\makeatletter
\newlabel{sites-cohomology-item-A-and-P}{{0EZ7}{source}{}{sourcepointer-7}{}}
\newlabel{sites-cohomology-item-A-and-P@cref}{{[originalreference][0][]0EZ7}{[1][0][]source}{}{}{}}
\makeatother
\crefname{originalreference}{original reference}{original references}
\Crefname{originalreference}{Original reference}{Original references}
\makeatletter
\newlabel{sites-cohomology-item-A-sheaf}{{0EZ6}{source}{}{sourcepointer-8}{}}
\newlabel{sites-cohomology-item-A-sheaf@cref}{{[originalreference][0][]0EZ6}{[1][0][]source}{}{}{}}
\makeatother
\crefname{originalreference}{original reference}{original references}
\Crefname{originalreference}{Original reference}{Original references}
\makeatletter
\newlabel{sites-cohomology-item-base-change-P}{{0EZ4}{source}{}{sourcepointer-9}{}}
\newlabel{sites-cohomology-item-base-change-P@cref}{{[originalreference][0][]0EZ4}{[1][0][]source}{}{}{}}
\makeatother
\crefname{originalreference}{original reference}{original references}
\Crefname{originalreference}{Original reference}{Original references}
\makeatletter
\newlabel{sites-cohomology-item-refine-tau-by-P}{{0EZ8}{source}{}{sourcepointer-10}{}}
\newlabel{sites-cohomology-item-refine-tau-by-P@cref}{{[originalreference][0][]0EZ8}{[1][0][]source}{}{}{}}
\makeatother
\crefname{originalreference}{original reference}{original references}
\Crefname{originalreference}{Original reference}{Original references}
\makeatletter
\newlabel{sites-cohomology-item-restriction-A}{{0EZ5}{source}{}{sourcepointer-11}{}}
\newlabel{sites-cohomology-item-restriction-A@cref}{{[originalreference][0][]0EZ5}{[1][0][]source}{}{}{}}
\makeatother
\crefname{originalreference}{original reference}{original references}
\Crefname{originalreference}{Original reference}{Original references}
\makeatletter
\newlabel{sites-cohomology-situation-compare}{{0EZ3}{source}{}{sourcepointer-12}{}}
\newlabel{sites-cohomology-situation-compare@cref}{{[originalreference][0][]0EZ3}{[1][0][]source}{}{}{}}
\makeatother
\crefname{originallemma}{original lemma}{original lemmas}
\Crefname{originallemma}{Original lemma}{Original lemmas}
\makeatletter
\newlabel{topologies-lemma-morphism-big-small-etale}{{021I}{source}{}{sourcepointer-13}{}}
\newlabel{topologies-lemma-morphism-big-small-etale@cref}{{[originallemma][0][]021I}{[1][0][]source}{}{}{}}
\makeatother
\title{The five standing properties hold for the comparison of the fppf topology with the etale topology, with finite morphisms as the class P}
\author{The Stacks Project Authors}
\date{stacks-project:etale-cohomology.tex:local-snapshot-20261011; GNU FDL 1.2 or later}
\begin{document}
\maketitle
\noindent\textit{Editorial note.} All delivered mathematical text is the author's own, copied from the pinned Stacks Project chapter file etale-cohomology.tex; no statement or proof step was rewritten, reordered or removed.\par\medskip
\noindent\textit{Editorial note.} The counted proof is the author's complete argument: the weak Serre verification through the four conditions, the four cited properties, and the appeal to the domination of an fppf covering by a finite morphism.\par\medskip
\noindent\textit{Editorial note.} The label lemma-compare-fppf-etale names the counted statement, and the proof environment immediately following it in the source is delivered verbatim as its proof.\par\medskip
\noindent\textit{Editorial note.} The counted statement is carried by the author's own lemma environment in the chapter source.\par\medskip
\noindent\textit{Editorial note.} Pre-open gate: python operations/agent-authoring/tools/precard\_gate.py --source etale-cohomology.tex --label lemma-compare-fppf-etale --start 21665 --end 21722 --local-id 3562 returned RESULT=OK (L3562 free, occupancy FREE, public interval clear). Candidate taken from lane-s4/halfB-priority.json.\par\medskip
\noindent\textit{Editorial note.} The statement refers to the five labelled items of the standing comparison situation of the sibling chapter; these are kept as external pointers with their canonical Stacks tags.\par\medskip
\noindent\textit{Editorial note.} The statement and proof are copied from the pinned chapter file.\par\medskip
\noindent\textit{Allowed standard prerequisites (editorial).}\begin{itemize}
\item Cohomological descent for the comparison of fppf with etale. Same chapter, the lemma establishing the properties of the comparison of the fppf and etale topologies, in particular the exactness and full faithfulness of the pullback and the exactness of the pushforward on short exact sequences of sheaves in its image. Supplies the weak Serre verification of the class of pulled back sheaves.
\item Characterization of weak Serre subcategories. Homology, the lemma listing the four conditions characterizing a weak Serre subcategory. Supplies the four conditions checked for the subcategory of sheaves pulled back from the etale site.
\item Base change and restriction for finite morphisms. Morphisms, the lemma that the base change of a finite morphism is finite, and Topologies, the lemma giving the commutative diagram of the morphism between the big and small etale sites. Supply property (item-base-change-P) and property (item-restriction-A) respectively.
\item The sheaf property and compatibility of the class with proper maps. Same chapter, the lemma describing the pullback along the projection for the fppf site, and part (4) of the lemma comparing the higher direct image for proper morphisms. Supply properties (item-A-sheaf) and (item-A-and-P).
\item Domination of an fppf covering by a finite morphism etale locally. More on Morphisms, the lemma that an fppf covering can be dominated, etale locally on the base, by a finite morphism. Supplies property (item-refine-tau-by-P), that is, the refinement of an fppf covering by a finite morphism.
\end{itemize}\medskip
\ifdefined\hypertarget\hypertarget{sourcepointer-1}{}\fi
\noindent\textit{Source reference.} Original lemma 0754: sources/stacks-chow/tags.txt, source lines 9118--9118. Prerequisite pointer invoked at proof line(s) 21687; it names a standard fact of a sibling chapter that the counted argument uses, and every step of the counted proof built on it is retained verbatim in the delivered body.\par\medskip
\ifdefined\hypertarget\hypertarget{sourcepointer-2}{}\fi
\noindent\textit{Source reference.} Original reference 0DDI: sources/stacks-etale-cohomology/pinned-original-source.tex, source lines 21308--21308. Prerequisite pointer invoked at proof line(s) 21689; it names a standard fact of the same chapter that the counted argument uses, and every step of the counted proof built on it is retained verbatim in the delivered body.\par\medskip
\ifdefined\hypertarget\hypertarget{sourcepointer-3}{}\fi
\noindent\textit{Source reference.} Original reference 0DDJ: sources/stacks-etale-cohomology/pinned-original-source.tex, source lines 21323--21323. Prerequisite pointer invoked at proof line(s) 21715; it names a standard fact of the same chapter that the counted argument uses, and every step of the counted proof built on it is retained verbatim in the delivered body.\par\medskip
\ifdefined\hypertarget\hypertarget{sourcepointer-4}{}\fi
\noindent\textit{Source reference.} Original reference 0DDL: sources/stacks-etale-cohomology/pinned-original-source.tex, source lines 21454--21454. Prerequisite pointer invoked at proof line(s) 21711; it names a standard fact of the same chapter that the counted argument uses, and every step of the counted proof built on it is retained verbatim in the delivered body.\par\medskip
\ifdefined\hypertarget\hypertarget{sourcepointer-5}{}\fi
\noindent\textit{Source reference.} Original lemma 0DET: sources/stacks-chow/tags.txt, source lines 16813--16813. Prerequisite pointer invoked at proof line(s) 21721; it names a standard fact of a sibling chapter that the counted argument uses, and every step of the counted proof built on it is retained verbatim in the delivered body.\par\medskip
\ifdefined\hypertarget\hypertarget{sourcepointer-6}{}\fi
\noindent\textit{Source reference.} Original lemma 01WL: sources/stacks-chow/tags.txt, source lines 2486--2486. Prerequisite pointer invoked at proof line(s) 21702; it names a standard fact of a sibling chapter that the counted argument uses, and every step of the counted proof built on it is retained verbatim in the delivered body.\par\medskip
\ifdefined\hypertarget\hypertarget{sourcepointer-7}{}\fi
\noindent\textit{Source reference.} Original reference 0EZ7: sources/stacks-chow/tags.txt, source lines 18717--18717. Prerequisite pointer invoked at proof line(s) 21678; it names a standard fact of a sibling chapter that the counted argument uses, and every step of the counted proof built on it is retained verbatim in the delivered body.\par\medskip
\ifdefined\hypertarget\hypertarget{sourcepointer-8}{}\fi
\noindent\textit{Source reference.} Original reference 0EZ6: sources/stacks-chow/tags.txt, source lines 18716--18716. Prerequisite pointer invoked at proof line(s) 21677; it names a standard fact of a sibling chapter that the counted argument uses, and every step of the counted proof built on it is retained verbatim in the delivered body.\par\medskip
\ifdefined\hypertarget\hypertarget{sourcepointer-9}{}\fi
\noindent\textit{Source reference.} Original reference 0EZ4: sources/stacks-chow/tags.txt, source lines 18714--18714. Prerequisite pointer invoked at proof line(s) 21675; it names a standard fact of a sibling chapter that the counted argument uses, and every step of the counted proof built on it is retained verbatim in the delivered body.\par\medskip
\ifdefined\hypertarget\hypertarget{sourcepointer-10}{}\fi
\noindent\textit{Source reference.} Original reference 0EZ8: sources/stacks-chow/tags.txt, source lines 18718--18718. Prerequisite pointer invoked at proof line(s) 21679; it names a standard fact of a sibling chapter that the counted argument uses, and every step of the counted proof built on it is retained verbatim in the delivered body.\par\medskip
\ifdefined\hypertarget\hypertarget{sourcepointer-11}{}\fi
\noindent\textit{Source reference.} Original reference 0EZ5: sources/stacks-chow/tags.txt, source lines 18715--18715. Prerequisite pointer invoked at proof line(s) 21676; it names a standard fact of a sibling chapter that the counted argument uses, and every step of the counted proof built on it is retained verbatim in the delivered body.\par\medskip
\ifdefined\hypertarget\hypertarget{sourcepointer-12}{}\fi
\noindent\textit{Source reference.} Original reference 0EZ3: sources/stacks-chow/tags.txt, source lines 18713--18713. Prerequisite pointer invoked at proof line(s) 21681; it names a standard fact of a sibling chapter that the counted argument uses, and every step of the counted proof built on it is retained verbatim in the delivered body.\par\medskip
\ifdefined\hypertarget\hypertarget{sourcepointer-13}{}\fi
\noindent\textit{Source reference.} Original lemma 021I: sources/stacks-chow/tags.txt, source lines 2671--2671. Prerequisite pointer invoked at proof line(s) 21707; it names a standard fact of a sibling chapter that the counted argument uses, and every step of the counted proof built on it is retained verbatim in the delivered body.\par\medskip
\begin{lemma}
\label{lemma-compare-fppf-etale}
Consider the comparison morphism
$\epsilon : (\Sch/S)_{fppf} \to (\Sch/S)_\etale$.
Let $\mathcal{P}$ denote the class of finite morphisms of schemes.
For $X$ in $(\Sch/S)_\etale$ denote
$\mathcal{A}'_X \subset \textit{Ab}((\Sch/X)_\etale)$
the full subcategory consisting of sheaves of the form
$\pi_X^{-1}\mathcal{F}$ with $\mathcal{F}$ in $\textit{Ab}(X_\etale)$.
Then Cohomology on Sites, Properties
(\ref{sites-cohomology-item-base-change-P}),
(\ref{sites-cohomology-item-restriction-A}),
(\ref{sites-cohomology-item-A-sheaf}),
(\ref{sites-cohomology-item-A-and-P}), and
(\ref{sites-cohomology-item-refine-tau-by-P})
of Cohomology on Sites, Situation
\ref{sites-cohomology-situation-compare} hold.
\end{lemma}

\begin{proof}
We first show that $\mathcal{A}'_X \subset \textit{Ab}((\Sch/X)_\etale)$
is a weak Serre subcategory by checking conditions (1), (2), (3), and (4)
of Homology, Lemma \ref{homology-lemma-characterize-weak-serre-subcategory}.
Parts (1), (2), (3) are immediate as $\pi_X^{-1}$ is exact and
fully faithful for example by Lemma \ref{lemma-cohomological-descent-etale}. If
$0 \to \pi_X^{-1}\mathcal{F} \to \mathcal{G} \to \pi_X^{-1}\mathcal{F}' \to 0$
is a short exact sequence in $\textit{Ab}((\Sch/X)_\etale)$
then $0 \to \mathcal{F} \to \pi_{X, *}\mathcal{G} \to \mathcal{F}' \to 0$
is exact by Lemma \ref{lemma-cohomological-descent-etale}.
Hence $\mathcal{G} = \pi_X^{-1}\pi_{X, *}\mathcal{G}$ is in
$\mathcal{A}'_X$ which checks the final condition.

\medskip\noindent
Cohomology on Sites, Property (\ref{sites-cohomology-item-base-change-P}) holds
by the existence of fibre products of schemes
and the fact that the base change of a finite morphism of
schemes is a finite morphism of schemes, see
Morphisms, Lemma \ref{morphisms-lemma-base-change-finite}.

\medskip\noindent
Cohomology on Sites, Property (\ref{sites-cohomology-item-restriction-A})
follows from the commutative diagram (3) in
Topologies, Lemma \ref{topologies-lemma-morphism-big-small-etale}.

\medskip\noindent
Cohomology on Sites, Property (\ref{sites-cohomology-item-A-sheaf}) is
Lemma \ref{lemma-describe-pullback-pi-fppf}.

\medskip\noindent
Cohomology on Sites, Property (\ref{sites-cohomology-item-A-and-P}) holds by
Lemma \ref{lemma-compare-higher-direct-image-proper} part (4).

\medskip\noindent
Cohomology on Sites, Property (\ref{sites-cohomology-item-refine-tau-by-P})
is implied by
More on Morphisms, Lemma
\ref{more-morphisms-lemma-dominate-fppf-etale-locally}.
\end{proof}

\end{document}
